\chaptersection{3}{Failure Mode and Effects Analysis}

La Failure Mode and Effects Analysis (FMEA) è una metodologia utilizzata
per identificare le possibili modalità di guasto di un sistema e valutarne
le conseguenze. L'analisi considera gli asset selezionati tra quelli
individuati nell'Asset Model, concentrandosi sugli elementi maggiormente
significativi per il funzionamento della pipeline di DermaTriage. Per
ciascun asset vengono esaminate le funzioni svolte, i relativi prerequisiti
e i possibili failure mode, individuandone le cause e gli effetti che
possono propagarsi agli altri componenti del sistema.

I failure mode considerati possono derivare sia da condizioni accidentali
sia da azioni intenzionali in grado di alterare il normale comportamento
dell'asset. L'eventuale corrispondenza con specifiche minacce di sicurezza
viene analizzata successivamente mediante STRIDE-AI.


% ============================================================
% A01
% ============================================================

\subsection{A01 -- Immagine dermatologica}

L'asset A01 rappresenta l'immagine della lesione dermatologica acquisita
per il caso clinico e utilizzata come principale informazione visiva
della pipeline. L'immagine può essere recuperata nell'ambito del consulto
gestito tramite B4 oppure fornita direttamente attraverso l'endpoint
\texttt{/analyze}.

All'interno della pipeline, l'immagine viene utilizzata nello Stage 1
da EfficientNet-B4 per la classificazione iniziale dell'urgenza e nello
Stage 2 da Qwen2-VL per la generazione della descrizione clinica della
lesione. La correttezza e l'adeguatezza dell'immagine costituiscono quindi
un prerequisito per le successive elaborazioni basate sull'informazione
visiva.


\subsubsection{Funzioni e prerequisiti}

Le principali funzioni dell'asset e i prerequisiti necessari al loro
svolgimento sono riportati nella Tabella~\ref{tab:a01-functions}.

\begin{table}[H]
    \centering
    \small
    \setstretch{1}
    \setlength{\tabcolsep}{4pt}
    \renewcommand{\arraystretch}{1.08}

    \begin{tabularx}{\textwidth}{
        >{\raggedright\arraybackslash}p{1.1cm}
        >{\raggedright\arraybackslash}p{3.8cm}
        X
    }
        \toprule

        \textbf{ID} &
        \textbf{Funzione} &
        \textbf{Prerequisiti} \\

        \midrule

        F01.1 &
        Rappresentazione visiva della lesione &
        Acquisizione di un'immagine che rappresenti in modo adeguato
        la lesione dermatologica oggetto del consulto. \\

        \midrule

        F01.2 &
        Associazione dell'immagine al caso clinico &
        Corretta corrispondenza tra l'immagine acquisita e il consulto
        o la richiesta di analisi a cui essa appartiene. \\

        \midrule

        F01.3 &
        Fornitura dell'input visivo alla pipeline di inferenza &
        Disponibilità dell'immagine in un formato utilizzabile e corretta
        trasmissione dell'input a EfficientNet-B4 nello Stage 1 e a Qwen2-VL
        nello Stage 2. \\

        \bottomrule
    \end{tabularx}

    \caption{Funzioni e prerequisiti dell'asset A01 -- Immagine dermatologica.}
    \label{tab:a01-functions}
\end{table}


\subsubsection{Failure mode}

Sulla base delle funzioni individuate vengono considerate le principali
modalità con cui l'asset A01 può non svolgere correttamente il proprio
compito, ovvero i failure mode.

\begin{table}[H]
    \centering
    \footnotesize
    \setstretch{1}
    \setlength{\tabcolsep}{3pt}
    \renewcommand{\arraystretch}{1.1}

    \begin{tabularx}{\textwidth}{
        >{\raggedright\arraybackslash}p{1.1cm}
        >{\raggedright\arraybackslash}p{3cm}
        >{\raggedright\arraybackslash}p{5cm}
        X
    }
        \toprule

        \textbf{ID} &
        \textbf{Failure mode} &
        \textbf{Possibili cause} &
        \textbf{Effetti} \\

        \midrule

        FM01.1 &
        Immagine di qualità insufficiente &
        Sfocatura, illuminazione inadeguata, risoluzione insufficiente,
        presenza di artefatti oppure condizioni di acquisizione che
        rendono poco distinguibili le caratteristiche della lesione. &
        L'immagine non rappresenta con sufficiente chiarezza le
        caratteristiche della lesione, riducendo l'affidabilità delle
        elaborazioni di EfficientNet-B4 e Qwen2-VL. \\

        \midrule

        FM01.2 &
        Lesione rappresentata in modo incompleto o non adeguato &
        Inquadratura parziale, porzione rilevante della lesione esclusa
        dall'immagine oppure acquisizione di una regione cutanea che non
        consente di osservare adeguatamente la lesione di interesse. &
        I modelli ricevono una rappresentazione visiva incompleta e possono
        produrre una classificazione o una descrizione clinica non
        rappresentative della lesione effettiva. \\

        \midrule

        FM01.3 &
        Immagine associata al caso clinico errato &
        Errata associazione tra immagine e consulto, selezione del
        documento non corrispondente al caso oppure errore nella gestione
        del riferimento utilizzato per recuperare l'immagine. &
        La pipeline elabora un'immagine non coerente con le informazioni
        cliniche della richiesta, producendo inferenze riferite al caso
        errato. \\

        \midrule

        FM01.4 &
        Manipolazione intenzionale dell'immagine &
        Alterazione deliberata del contenuto visivo, introduzione di
        perturbazioni o modifica dell'immagine prima dell'invio alla
        pipeline con l'obiettivo di influenzare l'elaborazione dei modelli. &
        EfficientNet-B4 e Qwen2-VL ricevono un input apparentemente valido
        ma alterato, con possibile modifica della classificazione
        dell'urgenza o della descrizione clinica prodotta. \\

        \midrule

        FM01.5 &
        Immagine non disponibile o non utilizzabile dalla pipeline &
        Formato non supportato, file danneggiato, errore nella lettura,
        mancato recupero dal servizio di origine oppure errore durante
        il trasferimento. &
        EfficientNet-B4 e Qwen2-VL non dispongono di un input visivo
        utilizzabile e non possono completare normalmente le rispettive
        elaborazioni. \\

        \bottomrule
    \end{tabularx}

    \caption{Failure mode individuati per l'asset A01 -- Immagine dermatologica.}
    \label{tab:a01-failure-modes}
\end{table}


\subsection{A13 -- Feedback clinico validato}

L'asset A13 rappresenta il feedback clinico prodotto durante la revisione
medica dei risultati generati da DermaTriage. Il feedback consente di
confermare oppure correggere l'esito del sistema e viene associato al
relativo caso clinico.

Le correzioni validate possono essere utilizzate nei meccanismi di
aggiornamento della pipeline. In particolare, contribuiscono
all'evoluzione dei prompt utilizzati dai modelli generativi e, nei casi
di disaccordo con la classificazione precedente, possono concorrere
all'attivazione del retraining del classificatore EfficientNet-B4.

Il feedback clinico rappresenta quindi un collegamento diretto tra la
revisione medica e il ciclo di aggiornamento del sistema.


\subsubsection{Funzioni e prerequisiti}

Le principali funzioni dell'asset e i prerequisiti necessari al loro
svolgimento sono riportati nella Tabella~\ref{tab:a13-functions}.

\begin{table}[H]
    \centering
    \small
    \setstretch{1}
    \setlength{\tabcolsep}{4pt}
    \renewcommand{\arraystretch}{1.08}

    \begin{tabularx}{\textwidth}{
        >{\raggedright\arraybackslash}p{1.1cm}
        >{\raggedright\arraybackslash}p{3.8cm}
        X
    }
        \toprule

        \textbf{ID} &
        \textbf{Funzione} &
        \textbf{Prerequisiti} \\

        \midrule

        F13.1 &
        Rappresentazione dell'esito della revisione medica &
        Disponibilità del risultato diagnostico da revisionare e corretta
        associazione tra valutazione del medico e relativo caso clinico. \\

        \midrule

        F13.2 &
        Fornitura di esempi validati per l'aggiornamento dei prompt &
        Disponibilità di correzioni cliniche correttamente registrate e
        utilizzabili dal meccanismo di evoluzione dei prompt. \\

        \midrule

        F13.3 &
        Supporto al retraining del classificatore &
        Presenza di correzioni che evidenziano un disaccordo rispetto alla
        classificazione precedente e disponibilità dei relativi dati clinici
        necessari al processo di aggiornamento del modello. \\

        \bottomrule
    \end{tabularx}

    \caption{Funzioni e prerequisiti dell'asset A13 -- Feedback clinico validato.}
    \label{tab:a13-functions}
\end{table}


\subsubsection{Failure mode}

Sulla base delle funzioni individuate vengono considerate le principali
modalità con cui l'asset A13 può non svolgere correttamente il proprio
compito, ovvero i failure mode.

\begin{table}[H]
    \centering
    \footnotesize
    \setstretch{1}
    \setlength{\tabcolsep}{3pt}
    \renewcommand{\arraystretch}{1.1}

    \begin{tabularx}{\textwidth}{
        >{\raggedright\arraybackslash}p{1.1cm}
        >{\raggedright\arraybackslash}p{3cm}
        >{\raggedright\arraybackslash}p{5cm}
        X
    }
        \toprule

        \textbf{ID} &
        \textbf{Failure mode} &
        \textbf{Possibili cause} &
        \textbf{Effetti} \\

        \midrule

        FM13.1 &
        Feedback contenente istruzioni o contenuti avversariali &
        Inserimento intenzionale nel feedback di testo costruito per
        influenzare i successivi meccanismi di evoluzione dei prompt oppure
        utilizzo di correzioni contenenti istruzioni non pertinenti al
        contenuto clinico. &
        Gli esempi utilizzati per aggiornare i prompt possono incorporare
        contenuti avversariali e influenzare il comportamento successivo
        di Qwen2-VL o BioMistral. \\

        \midrule

        FM13.2 &
        Feedback associato al caso clinico errato &
        Collegamento non corretto tra correzione e diagnosi, utilizzo di
        un identificativo del consulto errato oppure errore nella
        registrazione della relazione tra feedback e caso. &
        Una correzione valida può essere applicata a un caso differente,
        introducendo informazioni incoerenti nel ciclo di aggiornamento. \\

        \midrule

        FM13.3 &
        Feedback clinico incompleto o non disponibile &
        Mancata registrazione della revisione, campi necessari non
        valorizzati oppure indisponibilità della correzione nel momento
        in cui deve essere utilizzata. &
        I processi di aggiornamento dispongono di informazioni validate
        insufficienti e possono non attivarsi o operare su una base
        incompleta. \\

        \midrule

        FM13.4 &
        Feedback clinico alterato o non autentico &
        Modifica non autorizzata di una correzione registrata, sostituzione
        dell'esito validato oppure inserimento di feedback falsi da parte
        di un soggetto non autorizzato. &
        Il sistema può utilizzare informazioni manipolate per
        l'aggiornamento dei prompt o per il retraining, propagando
        l'alterazione alle versioni successive della pipeline. \\

        \bottomrule
    \end{tabularx}

    \caption{Failure mode individuati per l'asset A13 -- Feedback clinico validato.}
    \label{tab:a13-failure-modes}
\end{table}

% ============================================================
% A16
% ============================================================

\subsection{A16 -- Modello EfficientNet-B4}

L'asset A16 riguarda il modello EfficientNet-B4 utilizzato per la
classificazione iniziale dell'urgenza dermatologica e i parametri appresi
durante l'addestramento, caricati dal checkpoint attivo. A partire
dall'immagine della lesione, il modello produce una classe di urgenza
tra HIGH, MEDIUM e LOW, insieme alla confidence della classificazione e
alle probabilità associate alle diverse classi.

Il risultato dello Stage 1 viene utilizzato dalle successive elaborazioni
della pipeline e contribuisce alla sintesi clinica prodotta da BioMistral
e alla successiva determinazione della priorità P1--P4.


\subsubsection{Funzioni e prerequisiti}

Le principali funzioni dell'asset e i prerequisiti necessari al loro
svolgimento sono riportati nella Tabella~\ref{tab:a16-functions}.

\begin{table}[H]
    \centering
    \small
    \setstretch{1}
    \setlength{\tabcolsep}{4pt}
    \renewcommand{\arraystretch}{1.08}

    \begin{tabularx}{\textwidth}{
        >{\raggedright\arraybackslash}p{1.1cm}
        >{\raggedright\arraybackslash}p{3.8cm}
        X
    }
        \toprule

        \textbf{ID} &
        \textbf{Funzione} &
        \textbf{Prerequisiti} \\

        \midrule

        F16.1 &
        Estrazione delle caratteristiche visive &
        Disponibilità di un'immagine dermatologica acquisita correttamente
        (A01) e caricamento del modello con i relativi parametri. \\

        \midrule

        F16.2 &
        Classificazione dell'urgenza in HIGH, MEDIUM o LOW &
        Corretta elaborazione dell'immagine e utilizzo dei parametri
        previsti per la versione attiva del classificatore. \\

        \midrule

        F16.3 &
        Produzione della confidence e delle probabilità associate &
        Completamento corretto dell'inferenza e disponibilità dell'output
        numerico prodotto dal classificatore. \\

        \midrule

        F16.4 &
        Utilizzo della versione prevista del modello &
        Disponibilità del checkpoint attivo e corretta configurazione
        dei riferimenti utilizzati per il caricamento del modello (A23). \\

        \bottomrule
    \end{tabularx}

    \caption{Funzioni e prerequisiti dell'asset A16 -- Modello EfficientNet-B4 e relativi pesi.}
    \label{tab:a16-functions}
\end{table}


\subsubsection{Failure mode}

Sulla base delle funzioni individuate vengono considerate le principali
modalità con cui l'asset A16 può non svolgere correttamente il proprio
compito, ovvero i failure mode.

\begin{table}[H]
    \centering
    \footnotesize
    \setstretch{1}
    \setlength{\tabcolsep}{3pt}
    \renewcommand{\arraystretch}{1.1}

    \begin{tabularx}{\textwidth}{
        >{\raggedright\arraybackslash}p{1.1cm}
        >{\raggedright\arraybackslash}p{3cm}
        >{\raggedright\arraybackslash}p{5cm}
        X
    }
        \toprule

        \textbf{ID} &
        \textbf{Failure mode} &
        \textbf{Possibili cause} &
        \textbf{Effetti} \\

        \midrule

        FM16.1 &
        Estrazione inadeguata delle caratteristiche visive &
        Qualità insufficiente dell'immagine, presenza di artefatti visivi
        oppure inquadratura non adeguata della lesione. &
        Le caratteristiche ricavate dall'immagine non rappresentano
        adeguatamente la lesione, compromettendo l'informazione utilizzata
        per determinare la classe di urgenza. \\

        \midrule

        FM16.2 &
        Classificazione errata dell'urgenza &
        Limiti di generalizzazione del modello, distribuzione dei casi
        operativi differente da quella osservata durante l'addestramento
        oppure insufficiente capacità discriminativa tra le classi. &
        Il modello può sovrastimare o sottostimare l'urgenza, propagando
        l'errore agli stadi successivi e contribuendo a una valutazione
        finale non appropriata. \\

        \midrule

        FM16.3 &
        Confidence o probabilità non coerenti con la classificazione &
        Calibrazione non adeguata delle probabilità oppure comportamento
        eccessivamente confidente su campioni incerti o poco
        rappresentativi. &
        Gli stadi successivi ricevono una stima non attendibile
        dell'incertezza associata alla classificazione. \\

        \midrule

        FM16.4 &
        Alterazione o sostituzione non autorizzata dei pesi o del checkpoint &
        Modifica deliberata dei parametri del modello, sostituzione del
        checkpoint attivo con una versione differente o compromessa,
        oppure caricamento di una versione non validata. &
        Il classificatore può produrre inferenze sistematicamente diverse
        da quelle previste pur continuando ad apparire operativo. \\

        \midrule

        FM16.5 &
        Output di inferenza incompleto o non disponibile &
        Errore durante l'esecuzione del modello, incompatibilità tra input,
        modello e parametri caricati oppure anomalia software che interrompe
        l'elaborazione prima del completamento. &
        Lo Stage 1 non fornisce classe, confidence o probabilità utilizzabili,
        impedendo agli stadi dipendenti di proseguire normalmente. \\

        \bottomrule
    \end{tabularx}

    \caption{Failure mode individuati per l'asset A16 -- Modello EfficientNet-B4 e relativi pesi.}
    \label{tab:a16-failure-modes}
\end{table}
\subsection{A20 -- Regole dell'Adaptation Layer}

L'asset A20 rappresenta l'insieme delle regole utilizzate
dall'Adaptation Layer per convertire il livello di urgenza prodotto
dalla pipeline nella scala clinica P1--P4 utilizzata dal sistema di
triage.

La conversione considera la classe di urgenza HIGH, MEDIUM o LOW e,
per i casi classificati come HIGH, anche il livello di confidence.
Un caso HIGH con confidence superiore a 0.85 viene associato alla
priorità P1, mentre gli altri casi HIGH vengono assegnati a P2.
Le classi MEDIUM e LOW vengono convertite rispettivamente in P3 e P4.

L'Adaptation Layer produce inoltre le informazioni associate alla
priorità finale, tra cui lo specialista indicato e il relativo SLA.
Le regole di mapping costituiscono quindi il passaggio conclusivo tra
il risultato interno della pipeline di inferenza e la rappresentazione
clinica restituita dal sistema.


\subsubsection{Funzioni e prerequisiti}

Le principali funzioni dell'asset e i prerequisiti necessari al loro
svolgimento sono riportati nella Tabella~\ref{tab:a20-functions}.

\begin{table}[H]
    \centering
    \small
    \setstretch{1}
    \setlength{\tabcolsep}{4pt}
    \renewcommand{\arraystretch}{1.08}

    \begin{tabularx}{\textwidth}{
        >{\raggedright\arraybackslash}p{1.1cm}
        >{\raggedright\arraybackslash}p{3.8cm}
        X
    }
        \toprule

        \textbf{ID} &
        \textbf{Funzione} &
        \textbf{Prerequisiti} \\

        \midrule

        F20.1 &
        Traduzione dell'output di BioMistral nel risultato operativo di triage &
        Disponibilità di un output strutturato di BioMistral contenente il livello
        finale di urgenza e la relativa confidence, insieme alla corretta
        applicazione delle regole che determinano priorità P1--P4,
        specialista e SLA. \\

        \midrule

        F20.2 &
        Associazione delle indicazioni operative alla priorità assegnata &
        Disponibilità della priorità P1--P4 determinata e corretta
        corrispondenza con lo specialista e lo SLA previsti per il relativo
        livello di presa in carico. \\

        \bottomrule
    \end{tabularx}

    \caption{Funzioni e prerequisiti dell'asset A20 -- Regole dell'Adaptation Layer.}
    \label{tab:a20-functions}
\end{table}

\subsubsection{Failure mode}

Sulla base delle funzioni individuate vengono considerate le principali
modalità con cui l'asset A20 può non svolgere correttamente il proprio
compito, ovvero i failure mode.

\begin{table}[H]
    \centering
    \footnotesize
    \setstretch{1}
    \setlength{\tabcolsep}{3pt}
    \renewcommand{\arraystretch}{1.1}

    \begin{tabularx}{\textwidth}{
        >{\raggedright\arraybackslash}p{1.1cm}
        >{\raggedright\arraybackslash}p{3cm}
        >{\raggedright\arraybackslash}p{5cm}
        X
    }
        \toprule

        \textbf{ID} &
        \textbf{Failure mode} &
        \textbf{Possibili cause} &
        \textbf{Effetti} \\

        \midrule

        FM20.1 &
        Interpretazione errata dell'output di BioMistral &
        Output strutturato incompleto o inatteso, errore nella lettura
        dei campi restituiti, utilizzo di un valore differente da
        \texttt{final\_urgency} oppure gestione non corretta della
        confidence associata. &
        L'Adaptation Layer utilizza informazioni diverse da quelle
        effettivamente prodotte da BioMistral e può determinare una
        priorità P1--P4 non coerente con il risultato dello Stage 4. \\

        \midrule

        FM20.2 &
        Determinazione errata della priorità clinica &
        Definizione non corretta della corrispondenza tra HIGH, MEDIUM,
        LOW e P1--P4 oppure applicazione errata della soglia di confidence
        utilizzata per distinguere i casi P1 e P2. &
        Un output correttamente interpretato può essere trasformato in
        una priorità clinica diversa da quella prevista dalle regole
        dell'Adaptation Layer. \\

        \midrule

        FM20.3 &
        Alterazione intenzionale delle regole di mapping &
        Modifica non autorizzata della corrispondenza tra livelli di
        urgenza e priorità, variazione deliberata della soglia di
        confidence oppure sostituzione della logica di mapping prevista. &
        Il sistema può continuare a produrre risultati formalmente validi
        assegnando però sistematicamente priorità differenti da quelle
        previste. \\

        \midrule


        FM20.4 &
        Produzione incompleta o incoerente del risultato operativo &
        Mancata valorizzazione di priorità, specialista o SLA oppure
        associazione non coerente tra questi elementi nell'output
        prodotto dall'Adaptation Layer. &
        Il risultato finale del triage può risultare incompleto oppure
        contenere indicazioni operative non coerenti con la priorità
        assegnata. \\

        \bottomrule
    \end{tabularx}

    \caption{Failure mode individuati per l'asset A20 -- Regole dell'Adaptation Layer.}
    \label{tab:a20-failure-modes}
\end{table}

% ============================================================
% A21
% ============================================================

\subsection{A21 -- Piattaforma B4}

L'asset A21 rappresenta la piattaforma esterna B4, utilizzata come
principale punto di interazione tra DermaTriage e il flusso clinico.
Attraverso le API messe a disposizione dalla piattaforma, DermaTriage
può recuperare i consulti e i relativi documenti, acquisire l'immagine
dermatologica necessaria all'elaborazione e restituire il risultato
diagnostico prodotto dalla pipeline.

B4 supporta inoltre la gestione delle validazioni mediche associate ai
risultati di triage. Gli esiti clinici validati possono essere recuperati
da DermaTriage e utilizzati nei meccanismi di aggiornamento del sistema,
tra cui l'evoluzione dei prompt e il retraining del classificatore
EfficientNet-B4. L'accesso alle API B4 avviene mediante autenticazione
basata su Bearer JWT.


\subsubsection{Funzioni e prerequisiti}

Le principali funzioni dell'asset e i prerequisiti necessari al loro
svolgimento sono riportati nella Tabella~\ref{tab:a21-functions}.

\begin{table}[H]
    \centering
    \small
    \setstretch{1}
    \setlength{\tabcolsep}{4pt}
    \renewcommand{\arraystretch}{1.08}

    \begin{tabularx}{\textwidth}{
        >{\raggedright\arraybackslash}p{1.1cm}
        >{\raggedright\arraybackslash}p{3.8cm}
        X
    }
        \toprule

        \textbf{ID} &
        \textbf{Funzione} &
        \textbf{Prerequisiti} \\

        \midrule

        F21.1 &
        Raccolta e organizzazione delle informazioni del consulto &
        Presenza di una richiesta di triage correttamente identificata e
        associazione coerente tra dati clinici, informazioni demografiche
        e documenti caricati dal paziente. \\

        \midrule

        F21.2 &
        Messa a disposizione dei dati necessari a DermaTriage &
        Presenza sulla piattaforma delle informazioni e dei documenti
        associati al consulto e disponibilità delle API necessarie al loro
        recupero. \\

        \midrule

        F21.3 &
        Conservazione del risultato diagnostico e della validazione medica &
        Corretta identificazione del consulto a cui associare il risultato
        prodotto da DermaTriage e l'eventuale successiva valutazione del
        personale medico. \\

        \midrule

        F21.4 &
        Supporto al ciclo di aggiornamento di DermaTriage &
        Disponibilità di consulti, immagini ed esiti clinici validati
        utilizzabili dai processi di aggiornamento dei componenti della
        pipeline. \\

        \bottomrule
    \end{tabularx}

    \caption{Funzioni e prerequisiti dell'asset A21 -- Piattaforma B4.}
    \label{tab:a21-functions}
\end{table}


\subsubsection{Failure mode}

Sulla base delle funzioni individuate vengono considerate le principali
modalità con cui l'asset A21 può non svolgere correttamente il proprio
compito, ovvero i failure mode.

\begin{table}[H]
    \centering
    \footnotesize
    \setstretch{1}
    \setlength{\tabcolsep}{3pt}
    \renewcommand{\arraystretch}{1.1}

    \begin{tabularx}{\textwidth}{
        >{\raggedright\arraybackslash}p{1.1cm}
        >{\raggedright\arraybackslash}p{3cm}
        >{\raggedright\arraybackslash}p{5cm}
        X
    }
        \toprule

        \textbf{ID} &
        \textbf{Failure mode} &
        \textbf{Possibili cause} &
        \textbf{Effetti} \\

        \midrule

        FM21.1 &
        Dati del consulto incompleti o incoerenti &
        Informazioni cliniche mancanti, documenti non presenti,
        inconsistenza tra dati demografici, sintomi e documentazione
        oppure registrazione incompleta del consulto. &
        DermaTriage riceve una rappresentazione incompleta o incoerente
        del caso, riducendo l'affidabilità delle successive elaborazioni. \\

        \midrule

        FM21.2 &
        Manipolazione intenzionale dei dati o dei documenti del consulto &
        Modifica non autorizzata di informazioni cliniche, sostituzione
        dei documenti associati al consulto oppure alterazione dei
        riferimenti utilizzati per identificare il caso. &
        DermaTriage può elaborare informazioni intenzionalmente alterate
        o appartenenti a un caso differente, producendo un risultato di
        triage non corrispondente al caso reale. \\

        \midrule

        FM21.3 &
        Stato delle informazioni non correttamente aggiornato sulla piattaforma &
        Mancato aggiornamento del risultato diagnostico o della validazione,
        registrazione incompleta dello stato del consulto oppure persistenza
        di informazioni obsolete rispetto all'ultima elaborazione disponibile. &
        B4 può mantenere uno stato del caso non aggiornato, compromettendo
        la revisione medica o il successivo utilizzo del feedback clinico. \\

        \midrule

        FM21.4 &
        Servizi B4 non disponibili o non accessibili &
        Indisponibilità della piattaforma o delle API, problemi di
        autenticazione, mancato rilascio o rinnovo del Bearer JWT,
        interruzioni della comunicazione oppure sovraccarico intenzionale
        delle funzionalità esposte. &
        DermaTriage non può accedere alle funzionalità B4 necessarie,
        causando l'interruzione parziale o completa del flusso integrato. \\

        \bottomrule
    \end{tabularx}

    \caption{Failure mode individuati per l'asset A21 -- Piattaforma B4.}
    \label{tab:a21-failure-modes}
\end{table}

